\chapter{1958年河北卷}

\section{解答题}

\begin{problem}
求多项式 \(x^{376}-2x^{223}+5x^{145}-3x^{28}+3\) 用 \(x+1\) 除所得的余数.
\end{problem}

\begin{solution}
令 \(f(x)=x^{376}-2x^{223}+5x^{145}-3x^{28}+3\). 由余式定理，用 \(x+1\) 除 \(f(x)\) 所得的余数为 \(f(-1)\). 因为偶次幂的 \((-1)\) 次方为 \(1\)，奇次幂的 \((-1)\) 次方为 \(-1\)，所以
\[
f(-1)=(-1)^{376}-2(-1)^{223}+5(-1)^{145}-3(-1)^{28}+3=1+2-5-3+3=-2
\]
因此余数为 \(-2\).
\end{solution}

\begin{problem}
试求函数 \(y=\sin 2x\) 的周期，并作出它在一个周期内的图象.
\end{problem}

\begin{solution}
函数 \(y=\sin 2x\) 的角频率为 \(\omega=2\)，因此周期为 \(T=\frac{2\pi}{\lvert\omega\rvert}=\pi\). 取 \([0,\pi]\) 为一个周期，在该区间内依次计算
\(y(0)=0\)，\(y\left(\frac{\pi}{4}\right)=1\)，\(y\left(\frac{\pi}{2}\right)=0\)，\(y\left(\frac{3\pi}{4}\right)=-1\)，\(y(\pi)=0\). 将这五个点按正弦曲线顺次连接，即得一个周期内的图象.

\solutionfigure{\bitmapfigure[max width=0.42\linewidth]{img/1958/hebei/q2_fig1.png}}
\end{solution}

\begin{problem}
求 \(\left( \frac{1}{\sqrt[3]{x}}+\sqrt{x} \right)^{7}\) 的展开式中含 \(x\) 的一次幂的项.
\end{problem}

\begin{solution}
由二项式定理，记 \(k=0\)，\(1\)，\(\ldots\)，\(7\)，展开式的第 \(k+1\) 项为
\[
T_{k+1}=\mathrm C_7^k\left(x^{-\frac{1}{3}}\right)^{7-k}\left(x^{\frac{1}{2}}\right)^k=\mathrm C_7^k x^{-\frac{1}{3}(7-k)+\frac{k}{2}}
\]
要使该项含有 \(x\) 的一次幂，只需令幂指数等于 \(1\)，即
\[
-\frac{1}{3}(7-k)+\frac{k}{2}=1
\]
两边乘以 \(6\)，得 \(-2(7-k)+3k=6\)，所以 \(5k=20\)，从而 \(k=4\). 于是相应项的系数为
\[
\mathrm C_7^4=\mathrm C_7^3=\frac{7\cdot6\cdot5}{1\cdot2\cdot3}=35
\]
因此含 \(x\) 的一次幂的项为 \(35x\).
\end{solution}

\begin{problem}
设一圆锥的母线长为 \(l\)，底面半径为 \(r\). 用平行于底面的一个平面在这圆锥上截下一个小圆锥. 若小圆锥的侧面积等于原来圆锥的侧面积的一半，试求小圆锥的母线的长.

\bitmapfigure[float=inline,max width=0.40\linewidth]{img/1958/hebei/q4_fig1.png}
\end{problem}

\begin{solution}
设小圆锥的底面半径和母线长分别为 \(r'\) 和 \(l'\). 原圆锥的侧面积为 \(\pi rl\)，小圆锥的侧面积为 \(\pi r'l'\). 由截面平行于底面，小圆锥与原圆锥相似，所以
\[
\frac{r'}{r}=\frac{l'}{l}
\]
又由题意，
\[
\pi r'l'=\frac{1}{2}\pi rl
\]
由相似关系得 \(r'=\frac{rl'}{l}\)，代入面积关系并约去正数 \(\pi r\)，得 \(\frac{l'^2}{l}=\frac{l}{2}\)，即 \(l'^2=\frac{l^2}{2}\). 因为母线长为正数，所以
\[
l'=\frac{\sqrt{2}}{2}l
\]
\end{solution}

\begin{problem}
设 \(O\) 为 \(\triangle ABC\) 内的一点，线段 \(AO\) 的延长线交 \(BC\) 于 \(A_{1}\) 点，线段 \(BO\) 的延长线交 \(AC\) 于 \(B_{1}\) 点，线段 \(CO\) 的延长线交 \(AB\) 于 \(C_{1}\) 点，求证：

\begin{enumerate}
\item \(\frac{\triangle OBC}{\triangle ABC}=\frac{OA_{1}}{AA_{1}}\).
\item \(\frac{OA_{1}}{AA_{1}}+\frac{OB_{1}}{BB_{1}}+\frac{OC_{1}}{CC_{1}}=1\).
\end{enumerate}

\bitmapfigure[float=inline,max width=0.38\linewidth]{img/1958/hebei/q5_fig1.png}
\end{problem}

\begin{solution}
\begin{enumerate}
\item 设从 \(O\)、\(A\) 分别向 \(BC\) 作垂线，垂足为 \(D\)、\(E\)，对应的高分别为 \(h_1\)、\(h\). 因为两个三角形以同一条 \(BC\) 为底，所以
\[
\frac{\triangle OBC}{\triangle ABC}=\frac{\frac{1}{2}BC\cdot h_1}{\frac{1}{2}BC\cdot h}=\frac{h_1}{h}
\]
又因为 \(\triangle DOA_1\) 与 \(\triangle EAA_1\) 都是直角三角形，且在 \(A_1\) 处有一个锐角相等，所以两三角形相似，从而 \(\frac{h_1}{h}=\frac{OA_1}{AA_1}\). 因此 \(\frac{\triangle OBC}{\triangle ABC}=\frac{OA_1}{AA_1}\).

\item 设从 \(O\)、\(B\) 分别向 \(AC\) 作垂线，垂足为 \(F\)、\(G\)，对应的高分别为 \(h_2\)、\(h_B\). 以 \(AC\) 为底可得
\[
\frac{\triangle OCA}{\triangle ABC}=\frac{h_2}{h_B}
\]
由于点 \(F\)、\(O\)、\(B_1\) 和点 \(G\)、\(B\)、\(B_1\) 分别在两条垂线及同一直线 \(BB_1\) 上，直角三角形 \(\triangle FOB_1\) 与 \(\triangle GBB_1\) 相似，所以 \(\frac{h_2}{h_B}=\frac{OB_1}{BB_1}\)，从而 \(\frac{\triangle OCA}{\triangle ABC}=\frac{OB_1}{BB_1}\). 再设从 \(O\)、\(C\) 分别向 \(AB\) 作垂线，垂足为 \(H\)、\(I\)，对应的高分别为 \(h_3\)、\(h_C\). 以 \(AB\) 为底可得
\[
\frac{\triangle OAB}{\triangle ABC}=\frac{h_3}{h_C}
\]
直角三角形 \(\triangle HOC_1\) 与 \(\triangle ICC_1\) 相似，因此 \(\frac{h_3}{h_C}=\frac{OC_1}{CC_1}\)，从而 \(\frac{\triangle OAB}{\triangle ABC}=\frac{OC_1}{CC_1}\).
由于 \(O\) 在三角形内部，三角形 \(OBC\)、\(OCA\)、\(OAB\) 恰好把 \(\triangle ABC\) 分成三个部分，所以
\[
\begin{aligned}
\frac{OA_1}{AA_1}+\frac{OB_1}{BB_1}+\frac{OC_1}{CC_1}
&=\frac{\triangle OBC+\triangle OCA+\triangle OAB}{\triangle ABC}\\
&=\frac{\triangle ABC}{\triangle ABC}=1.
\end{aligned}
\]
\end{enumerate}
\end{solution}

\begin{problem}
解方程组 \(\begin{cases}
\cos 2x+\cos 2y=\frac{\sqrt{3}}{2}, &\circled{1}\\
x+y=\arccos\left( \frac{1}{2} \right) &\circled{2}
\end{cases}\)
\end{problem}

\begin{solution}
由反余弦函数的定义，\(\arccos\left(\frac{1}{2}\right)=\frac{\pi}{3}\)，所以第二个方程给出 \(x+y=\frac{\pi}{3}\). 利用和差化积公式，第一式左端为
\[
\cos 2x+\cos 2y=2\cos(x+y)\cos(x-y)=2\cdot\frac{1}{2}\cos(x-y)=\cos(x-y)
\]
因此 \(\cos(x-y)=\frac{\sqrt{3}}{2}\)，从而
\(x-y=2k\pi\pm\frac{\pi}{6}\)，\(k\in\Z\).
当 \(x-y=2k\pi+\frac{\pi}{6}\) 时，联立 \(x+y=\frac{\pi}{3}\) 得 \(x=\frac{\pi}{4}+k\pi\)，\(y=\frac{\pi}{12}-k\pi\). 当 \(x-y=2k\pi-\frac{\pi}{6}\) 时，联立得 \(x=\frac{\pi}{12}+k\pi\)，\(y=\frac{\pi}{4}-k\pi\). 这两族解均满足原方程组，因此方程组的解为
\(\left(x,y\right)=\left(\frac{\pi}{4}+k\pi,\frac{\pi}{12}-k\pi\right)\quad\text{或}\quad \left(\frac{\pi}{12}+k\pi,\frac{\pi}{4}-k\pi\right)\)，\(k\in\Z\).
\end{solution}

\begin{problem}
一农场有甲，乙两台打麦机，甲机的工作效率是乙机的 \(2\) 倍. 先只用甲机打完全部麦子的 \(\frac{2}{3}\)，然后只用乙机继续打完，所需要的时间比同时用两台打麦机打完全部麦子所需要的时间多 \(4\) 天，问分别用一台打麦机打完所有的麦子各需要多少天？
\end{problem}

\begin{solution}
设全部麦子为 \(1\)，乙机每天的工作量为 \(x\)，则 \(x>0\)，甲机每天的工作量为 \(2x\). 先用甲机完成 \(\frac{2}{3}\) 的工作量，再用乙机完成 \(\frac{1}{3}\) 的工作量，所需时间为 \(\frac{2/3}{2x}+\frac{1/3}{x}\). 两机同时工作时每天完成 \(3x\)，完成全部工作量所需时间为 \(\frac{1}{3x}\). 根据题意，
\[
\frac{2/3}{2x}+\frac{1/3}{x}=\frac{1}{3x}+4
\]
左端化为 \(\frac{2}{3x}\)，所以 \(\frac{2}{3x}=\frac{1}{3x}+4\)，即 \(\frac{1}{3x}=4\)，解得 \(x=\frac{1}{12}\). 因此，单用甲机完成全部麦子需要 \(\frac{1}{2x}=6\) 天，单用乙机完成全部麦子需要 \(\frac{1}{x}=12\) 天.
\end{solution}

\begin{problem}
已知方程 \(4x^{2}-(4\cos\theta)x+1=0\) 有实根，求 \(\theta\) 的一切可能的值.
\end{problem}

\begin{solution}
这是关于 \(x\) 的一元二次方程，且二次项系数 \(4\neq0\). 方程有实根的必要且充分条件是判别式不小于 \(0\)，即
\[
\Delta=(-4\cos\theta)^2-4\cdot4\cdot1=16\cos^2\theta-16\geqslant0
\]
所以 \(\cos^2\theta\geqslant1\). 又因为对任意实数 \(\theta\) 都有 \(\lvert\cos\theta\rvert\leqslant1\)，故只能是 \(\cos^2\theta=1\)，即 \(\cos\theta=\pm1\). 因而 \(\theta=k\pi\)，其中 \(k\in\Z\). 反过来，当 \(\theta=k\pi\) 时 \(\Delta=0\)，方程确有实根，所以全部可能值为 \(\theta=k\pi\)，\(k\in\Z\).
\end{solution}
